% Generated draft: requires mathematical and source review.
\documentclass[11pt]{article}
\usepackage[margin=25mm]{geometry}
\usepackage{fontspec}
\setmainfont{Latin Modern Roman}
\usepackage{amsmath,amssymb,mathtools,unicode-math}
\setmathfont{Latin Modern Math}
\usepackage{enumitem,needspace,xurl,hyperref}
\setlength{\parindent}{0pt}
\setlength{\parskip}{5pt}
\setlength{\emergencystretch}{3em}
\newcommand{\N}{\mathbb{N}}
\newcommand{\Z}{\mathbb{Z}}
\newcommand{\Q}{\mathbb{Q}}
\newcommand{\R}{\mathbb{R}}
\newcommand{\C}{\mathbb{C}}
\newcommand{\F}{\mathbb{F}}
\newcommand{\A}{\mathbb{A}}
\newcommand{\PP}{\mathbb{P}}
\newcommand{\E}{\mathbb{E}}
\newcommand{\eps}{\varepsilon}
\newcommand{\dd}{\,\mathrm{d}}
\begin{document}
\section*{UnsolvedMath ID: 20002895}
\textbf{Source problem code:} catalogue entry
\subsection{Definitions and mathematical statement}
Let $M\subset\C^N$ be a compact real-analytic CR manifold of hypersurface type, of real dimension $2d-1$ and CR dimension $d-1$, whose CR structure is spherical: each point is locally CR equivalent to a piece of the unit sphere in $\C^d$. Suppose $M$ bounds a complex analytic variety $V$. Huang's question is to describe the possible singularities of $V$. Boundedness, purity of dimension, normality, and the precise boundary/filling convention are not fully specified in the extracted problem, and cannot be imposed silently on its general classification target. The source reports that when $M$ is algebraic, $V$ has ``at most one isolated singular point'', citing [HJ], and also points to survey [H]. The precise theorem hypotheses and whether the quoted phrase excludes all nonisolated singularities require the missing references.

The supplied correction of the source example is
\[
f:\C^2\to\C^3,\quad f(z_1,z_2)=(z_1^2,\sqrt2z_1z_2,z_2^2),\qquad
C=\{Q(w)=w_2^2-2w_1w_3=0\}.
\]
The full affine image $f(\C^2)$ is $C$, while the ball image and its boundary are
\[
V=f(\mathbb B^2)=C\cap\mathbb B^3,\qquad M=f(S^3)=C\cap S^5,
\]
since $\|f(z)\|^2=\|z\|^4$. The gradient $\nabla Q=(-2w_3,2w_2,-2w_1)$ vanishes only at $0$. That singularity belongs to $C$ and $V$, whereas $M$ is smooth and avoids $0$. Thus the source's phrase ``$M$ has exactly one singularity'' must be corrected to the filling, and its ball image must not be identified with the entire affine cone. The map identifies $z$ with $-z$ and is a generically two-to-one quotient, not a normalization map of the already normal cone.

More generally the aid gives the following explicit family. For $d\geq2$, $k\geq2$, put $L=\binom{d+k-1}{k}$ and
\[
\nu_k:\C^d\to\C^L,\qquad
\nu_k(z)=\bigl(\sqrt{k!/\alpha!}\,z^\alpha\bigr)_{|\alpha|=k},\qquad
\|\nu_k(z)\|^2=\|z\|^{2k}.
\]
Let $\mu_k$ act by scalar multiplication on $\C^d$. Its image $C_{d,k}=\nu_k(\C^d)$ realizes the normal affine quotient $\C^d/\mu_k$, whose unique singularity is the vertex, of cyclic quotient type $\frac1k(1,\ldots,1)$. Set
\[
V_{d,k}=\nu_k(\mathbb B^d)=C_{d,k}\cap\mathbb B^L,\qquad
M_{d,k}=\nu_k(S^{2d-1})=C_{d,k}\cap S^{2L-1}.
\]
The boundary is a smooth real-algebraic spherical CR lens space $S^{2d-1}/\mu_k$, with $\pi_1(M_{d,k})\cong\Z/k\Z$. Thus it records $k$ through its fundamental group and bounds an explicit filling with one isolated cyclic quotient singularity. The parametrization is generically $k$-to-one, rather than birational, and must not be called a normalization.

The aid also reports that known algebraic normal fillings are finite ball quotients, and that published Grauert-tube constructions give non-quotient examples in the real-analytic category. Those broad literature assertions come without complete citations or all their hypotheses and are not promoted here to a universal theorem for every possible $V$. The explicit quotient family gives examples, not a classification of all real-analytic spherical CR fillings.
\subsection{Short English statement}
The problem asks which singularities can occur in fillings of compact real-analytic spherical CR manifolds. The supplied Veronese family gives smooth lens-space boundaries whose fillings have one isolated cyclic quotient singularity. It also corrects two slips in the source example: the singular point lies in the filling, and the image of the ball is only the bounded part of the cone. The broad classification and cited general results remain qualified by missing hypotheses and references.
\subsection{Sources}
\begin{itemize}
\item Local registry: permanent catalogue entry, exact source code catalogue entry; \texttt{aim-several-complex-variables-notes.json}, \texttt{notes[52]}.
\item AIM workshop list \emph{Complexity of mappings in CR geometry}, Problem 13, ``Singularities of varieties'', proposed by Xiaojun Huang, \url{https://aimath.org/WWN/complexitycrmap/complexitycrmap.pdf}.
\item Source reference labels [HJ] for the algebraic-boundary assertion and [H] for a survey. No bibliography entries, titles, authors beyond the proposer, years, or complete authoritative citations are included; the separate original literature field is empty.
\item Supplied research attempt 1, registry updated 14 August 2026: partial result, low novelty confidence, machine-generated and not peer reviewed. It supplies the corrected example and generalized Veronese family, and reports finite-ball-quotient and Grauert-tube literature assertions without complete citations.
\item The verification field identifies the problem and proposer but ends before reproducing further checked wording.
\end{itemize}
\end{document}
